\documentclass[12pt, a4paper]{article}

\usepackage[utf8]{inputenc}
\usepackage[T1]{fontenc}
\usepackage[brazilian]{babel}
\usepackage{geometry}
\usepackage{booktabs}
\usepackage{enumitem}
\usepackage{hyperref}
\usepackage{parskip}

\geometry{margin=2.2cm}

\hypersetup{colorlinks=true, urlcolor=blue, linkcolor=black}

\begin{document}

\begin{center}
    {\large\textbf{Proposta de Trabalho Final --- CMP570 Fotografia Computacional}}\\[0.3em]
    {\normalsize\textbf{Aluno:} Marcel Fernandes Gomes \quad
     \textbf{Data:} \today}
\end{center}

\hrule
\vspace{0.6em}

% -------------------------------------------------------------------
\section*{1. Título}

\textbf{Superresolução no Domínio de Frequências: Fine-tuning do SATLAS com Focal Frequency
Loss para Correção do Gap Espectral Planet--Aerofoto}

% -------------------------------------------------------------------
\section*{2. Descrição}

O modelo SATLAS~\cite{satlassuperres} realiza superresolução (SR) de imagens de satélite
Sentinel-2 para resolução aerofotogramétrica (NAIP), usando uma arquitetura ESRGAN com
entrada de múltiplas imagens (16 frames, 48 canais). Ao substituir Sentinel-2 por imagens
PlanetScope reamostradas para 10\,m/px surge um \emph{gap espectral}: o PlanetScope preserva
alta frequência que o Sentinel-2 não possui, fazendo o modelo SATLAS gerar alucinações de textura.

Este trabalho realiza fine-tuning do SATLAS com pares Planet $\rightarrow$ aerofoto locais
(Porto Alegre, $\sim$224\,km$^2$) e incorpora a \textbf{Focal Frequency Loss}
(FFL)~\cite{jiang2021} como loss auxiliar no domínio espectral. A FFL computa a DFT 2D da
saída do gerador e do ground truth e penaliza o erro componente a componente no espectro,
com um mecanismo focal que aumenta progressivamente o peso nas frequências onde o modelo
persiste errando. Um \emph{ablation study} com três variantes (baseline sem FFL,
$\lambda_{\text{FFL}}=0{,}1$ e $\lambda_{\text{FFL}}=1{,}0$) quantifica o impacto espectral
sem introduzir outros fatores de confusão. O projeto conecta diretamente aos tópicos de
Fourier 2D, aliasing e Teorema da Convolução.

% -------------------------------------------------------------------
\section*{3. Objetivos}

\begin{enumerate}[leftmargin=*, topsep=2pt, itemsep=1pt]
    \item Reduzir o gap espectral (métrica PSD\_L2) entre a SR do SATLAS e a aerofoto GT.
    \item Quantificar o trade-off entre fidelidade espectral (PSD\_L2) e qualidade de pixel
          (SSIM, cPSNR) para dois valores de $\lambda_{\text{FFL}}$.
    \item Demonstrar que a FFL é um complemento viável às losses pixel/perceptual/GAN já
          existentes no SATLAS, sem afetar a compatibilidade com treinos anteriores.
\end{enumerate}

% -------------------------------------------------------------------
\section*{4. Materiais e obtenção}

\begin{itemize}[leftmargin=*, topsep=2pt, itemsep=1pt]
    \item \textbf{Imagens PlanetScope (LR):} 16 GeoTIFFs RGB, cobertura $\sim$13,9\,$\times$\,16,1\,km,
          cenas mensais de jan/2025 a abr/2026 --- obtidas via Plataforma Brasil+.
    \item \textbf{Aerofoto ground truth (HR):} GeoTIFF Convênio do Estado do Rio Grande do Sul e Exército Brasileiro a $\sim$0,35\,m/px, disponíveis no IEDE-RS ou via 1º Centro de Geoinformação.
    \item \textbf{Modelo pré-treinado SATLAS:} checkpoint público \texttt{esrgan\_16S2.pth}
          (\url{https://github.com/allenai/satlas-super-resolution}).
    \item \textbf{Infraestrutura:} GPU Tesla V100-PCIE-32GB (servidor do PPGC/UFRGS,
          CUDA 12.9), Python 3.12, PyTorch $\geq$2.0, BasicSR.
    \item \textbf{Código FFL:} implementação própria baseada em Jiang et al.\ (2021),
          integrada ao SATLAS de forma retrocompatível.
\end{itemize}

% -------------------------------------------------------------------
\section*{5. Resultados esperados}

Espera-se que a FFL com $\lambda=1{,}0$ reduza o PSD\_L2 em pelo menos 5\% frente ao
baseline, confirmando que a penalização explícita no domínio espectral corrige o gap
Planet/aerofoto. O custo em SSIM deve ser pequeno ($|\Delta|<0{,}01$), caracterizando
um trade-off favorável para aplicações de fidelidade espectral em SR de imagens de satélite.
Resultados preliminares (Rodada~1, 1.476 chips) já indicaram redução de 7,2\% no PSD\_L2
com $\lambda=1{,}0$ e queda de apenas 0,006 no SSIM.

% -------------------------------------------------------------------
\section*{6. Etapas (\textit{milestones})}

\begin{enumerate}[leftmargin=*, topsep=2pt, itemsep=1pt]
    \item Preparação dos chips (LR/HR co-registrados, 9.360 pares) --- \textit{concluída}.
    \item Integração da FFL ao código SATLAS e configuração dos YAMLs do ablation --- \textit{concluída}.
    \item Treinamento das 3 variantes (20.000 iters cada, $\sim$45\,min/variante na V100) --- \textit{em andamento}.
    \item Avaliação quantitativa (SSIM, cPSNR, PSD\_L2) e geração das figuras comparativas.
    \item Redação do relatório final e preparação dos slides.
\end{enumerate}

% -------------------------------------------------------------------
\section*{7. Cronograma}

\begin{center}
\begin{tabular}{lll}
    \toprule
    Semana & Período & Atividade \\
    \midrule
    12 & 26/05 -- 01/06 & Treinamento Rodada 2 (em background no servidor) \\
    13 & 02/06 -- 08/06 & Avaliação, geração de figuras, análise dos resultados \\
    13--14 & 09/06 -- 15/06 & Redação do relatório final (LaTeX) \\
    14--15 & 16/06 -- 22/06 & Preparação e revisão dos slides; entrega final \\
    \bottomrule
\end{tabular}
\end{center}

\vspace{0.5em}
\textbf{Prazo final:} 22 de junho de 2026.

% -------------------------------------------------------------------
\begin{thebibliography}{9}
\bibitem{jiang2021}
    L.~Jiang, B.~Dai, W.~Wu, C.~C. Loy,
    ``Focal Frequency Loss for Image Reconstruction and Synthesis,''
    \textit{ICCV}, pp.~13919--13929, 2021.
\bibitem{satlassuperres}
    P.~Wolters, F.~Bastani, A.~Kembhavi,
    ``Zooming Out on Zooming In: Advancing Super-Resolution for Remote Sensing,''
    \textit{arXiv:2311.18082}, 2023.
    \url{https://arxiv.org/abs/2311.18082}.
\end{thebibliography}

\end{document}
